% © 2026 Ing. David Jaroš — CC BY-NC-ND 4.0
%
% This work is licensed under a Creative Commons Attribution-NonCommercial-NoDerivatives
% 4.0 International License (CC BY-NC-ND 4.0).
%
% See LICENSE.md for full license text.

% UBT-AI-PROVENANCE-BEGIN
% schema: ubt-ai-provenance/v1
% tier: C_working
% ai_assistance: disclosed
% human_review: risk-based
% editorial_responsibility: Ing. David Jaroš
% policy: ../../AI_PROVENANCE.md
% notice: Working material; exhaustive human review is not claimed.
% UBT-AI-PROVENANCE-END

\documentclass[12pt,a4paper]{article}
\usepackage[a4paper,margin=2.5cm]{geometry}
\usepackage{amsmath,amssymb,amsthm,tcolorbox}
\tcbuselibrary{skins,breakable}
\newtheorem{lemma}{Lemma}[section]
\newtheorem{theorem}[lemma]{Theorem}
\newtcolorbox{statusbox}[1][]{colback=gray!8!white,colframe=black!60,arc=3pt,boxrule=1pt,title=Status Summary,#1}

\title{SU(2) Scherk--Schwarz Twist and \(N_{\mathrm{eff}}=12\) in UBT}
\author{Ing.\ David Jaroš}
\date{2026-05-12}

\begin{document}
\maketitle

\section{Twisted boundary condition on \(S^1_\psi\)}
Define
\[
\Theta(\psi+2\pi R_\psi)=\sigma_3\,\Theta(\psi)\,\sigma_3,
\qquad
\sigma_3=\mathrm{diag}(+1,-1).
\]
This SU(2) twist replaces the earlier \(\tau\)-twist route.

\section{Symmetry of the kinetic action}
\begin{lemma}[Twist invariance of \(S_{\mathrm{kin}}\)]
For
\[
S_{\mathrm{kin}}=\int d^4x\;\mathrm{Tr}\!\left[(D_\mu\Theta)^\dagger(D^\mu\Theta)\right],
\]
and electromagnetic-covariant derivative satisfying
\(D_\mu(\sigma_3\Theta\sigma_3)=\sigma_3(D_\mu\Theta)\sigma_3\),
one has \(S_{\mathrm{kin}}\) invariant under
\(\Theta\mapsto\sigma_3\Theta\sigma_3\).
\end{lemma}

\begin{proof}
\[
\begin{aligned}
\mathrm{Tr}\!\left[(\sigma_3 D_\mu\Theta\sigma_3)^\dagger(\sigma_3 D^\mu\Theta\sigma_3)\right]
&=\mathrm{Tr}\!\left[\sigma_3(D_\mu\Theta)^\dagger\sigma_3\sigma_3(D^\mu\Theta)\sigma_3\right]\\
&=\mathrm{Tr}\!\left[\sigma_3(D_\mu\Theta)^\dagger(D^\mu\Theta)\sigma_3\right]\\
&=\mathrm{Tr}\!\left[(D_\mu\Theta)^\dagger(D^\mu\Theta)\right],
\end{aligned}
\]
using \(\sigma_3^2=1\) and cyclic invariance of trace.
\end{proof}

\section{KK spectrum under the SU(2) twist}
Write
\[
\Theta=\begin{pmatrix}a & b\\ c & d\end{pmatrix}.
\]
From
\(\Theta(\psi+2\pi R_\psi)=\sigma_3\Theta(\psi)\sigma_3\),
\[
\begin{pmatrix}
a(\psi+2\pi R) & b(\psi+2\pi R)\\
c(\psi+2\pi R) & d(\psi+2\pi R)
\end{pmatrix}
=
\begin{pmatrix}
a(\psi) & -b(\psi)\\
-c(\psi) & d(\psi)
\end{pmatrix}.
\]
Hence
\[
a,d:\ \text{periodic}\Rightarrow m_n=\frac{n}{R_\psi},
\qquad
b,c:\ \text{anti-periodic}\Rightarrow m_n=\frac{n+1/2}{R_\psi}.
\]
The charged off-diagonal fields are \(b,c\); diagonal \(a,d\) are neutral.

\section{Half-integer tower and fermionic effective weight}
For anti-periodic KK tower \((n+1/2)/R\), the partition-function sector is
\(\vartheta_2\) and the standard Scherk--Schwarz result gives Weyl-fermion-equivalent one-loop weight [STD].
Thus each charged complex field \(b\) and \(c\) contributes as one Weyl fermion in the effective photon vacuum polarization counting.

\section{Main counting theorem}
\begin{theorem}[\(N_{\mathrm{eff}}=12\) from SU(2) twist]
With three compact phase directions \((\psi,\chi,\xi)\), two charged off-diagonal fields \((b,c)\), and two Weyl helicity weights from the half-integer tower,
\[
N_{\mathrm{eff}}=N_{\mathrm{phases}}\times N_{\mathrm{charged}}\times N_{\mathrm{helicity}}
=3\times2\times2=12.
\]
Therefore
\[
B_0=\frac{2\pi}{3}N_{\mathrm{eff}}=\frac{2\pi}{3}\cdot12=8\pi.
\]
\end{theorem}

\begin{statusbox}
SU(2) twist symmetry \(S[\Theta]\): \textbf{[L0]}\\
Anti-per. BC pro \(b,c\): \textbf{[L1]}\\
Half-integer KK = Weyl fermion contribution: \textbf{[STD]}\\
\(N_{\mathrm{eff}} = 3\times2\times2 = 12\): \textbf{[L1]}\\
\(B_0 = 8\pi\): \textbf{[L1]}\\
Alpha je odvozena: \textbf{NE}
\end{statusbox}

\end{document}
